\section{Síntesis y ampliación}

\subsection{Repaso de los conocimientos clave}

Esta clase presentó de manera sistemática las dos arquitecturas fundamentales del modelado generativo profundo:

\textbf{1. Variational Autoencoder (VAE)}:
\begin{itemize}
    \item Comprime los datos en un espacio latente probabilístico mediante una arquitectura de codificador-decodificador
    \item Función de pérdida = pérdida de reconstrucción + regularización por divergencia KL
    \item El Reparameterization Trick permite un entrenamiento diferenciable de extremo a extremo
    \item El espacio latente tiene continuidad e integridad, lo que permite manipular características de manera interpretable
\end{itemize}

\textbf{2. Generative Adversarial Network (GAN)}:
\begin{itemize}
    \item Juego adversarial minimax entre el generador y el discriminador
    \item No necesita modelar la densidad de forma explícita: aprende directamente la transformación de la distribución
    \item La calidad de generación es excelente, pero la estabilidad del entrenamiento es el reto central
    \item Dio origen a variantes potentes como Progressive GAN, Conditional GAN y CycleGAN
\end{itemize}

\begin{importantbox}{Una perspectiva unificada de los modelos generativos}
Tanto el VAE como el GAN comparten el mismo objetivo: aprender la distribución de probabilidad subyacente de los datos. El VAE logra este objetivo mediante una estimación explícita de la densidad, mientras que el GAN aprende la distribución de los datos de manera implícita a través del entrenamiento adversarial. Ambos métodos tienen ventajas y desventajas propias, y representan dos filosofías fundamentales del modelado generativo: \textbf{modelado explícito} frente a \textbf{muestreo implícito}.
\end{importantbox}

\subsection{Del VAE/GAN a los modelos generativos actuales}

El VAE (2014) y el GAN (2014) presentados en esta clase son dos hitos en el campo de los modelos generativos. Después de ellos, este campo siguió avanzando con rapidez:

\begin{itemize}
    \item \textbf{Diffusion Models} (modelos de difusión, 2020--): generan datos mediante un proceso gradual de eliminación de ruido, y se han convertido en el método dominante para la generación de imágenes (por ejemplo, Stable Diffusion, DALL-E 2/3)
    \item \textbf{Flow-based Models} (modelos de flujo): modelan la densidad de forma explícita mediante transformaciones invertibles, como Glow y RealNVP
    \item \textbf{Autoregressive Models} (modelos autorregresivos): como la serie GPT, que modelan la distribución de probabilidad condicional token por token
    \item \textbf{Score-based Models}: generan muestras estimando el campo de gradiente de la distribución de los datos
\end{itemize}

\begin{knowledgebox}{¿Por qué los Diffusion Models terminaron superando a los demás?}
Los Diffusion Models combinan las ventajas del VAE y del GAN: al igual que el VAE, tienen un proceso de entrenamiento estable (basado en la optimización de la verosimilitud), y al igual que el GAN, son capaces de generar muestras de alta calidad. Mediante un proceso markoviano de eliminación gradual de ruido, los Diffusion Models evitan tanto la inestabilidad de entrenamiento del GAN como el problema de generación borrosa del VAE. Las clases posteriores de MIT 6.S191 explicarán los Diffusion Models con detalle.
\end{knowledgebox}

\subsection{Lecturas adicionales}

\begin{itemize}
    \item Kingma \& Welling, ``Auto-Encoding Variational Bayes,'' ICLR 2014 --- el artículo fundacional del VAE
    \item Goodfellow et al., ``Generative Adversarial Nets,'' NeurIPS 2014 --- el artículo fundacional del GAN
    \item Higgins et al., ``$\beta$-VAE: Learning Basic Visual Concepts with a Constrained Variational Framework,'' ICLR 2017
    \item Isola et al., ``Image-to-Image Translation with Conditional Adversarial Networks (Pix2Pix),'' CVPR 2017
    \item Karras et al., ``Progressive Growing of GANs,'' ICLR 2018
    \item Zhu et al., ``Unpaired Image-to-Image Translation using Cycle-Consistent Adversarial Networks (CycleGAN),'' ICCV 2017
    \item Brock et al., ``Large Scale GAN Training for High Fidelity Natural Image Synthesis (BigGAN),'' ICLR 2019
    \item Ho et al., ``Denoising Diffusion Probabilistic Models,'' NeurIPS 2020 --- Diffusion Models
    \item Página principal del curso MIT 6.S191: \href{https://introtodeeplearning.com}{introtodeeplearning.com}
    \item Código del Lab 2: \href{https://github.com/MITDeepLearning/introtodeeplearning}{github.com/MITDeepLearning/introtodeeplearning}
\end{itemize}

%% --- Fin del contenido --- %%

\end{document}
